\section{Conclusion}\label{sec:conclusion}

Measured against a provable optimum, energy-efficient WSN routing falls into clear levels. A recent multi-agent DRL router and a published energy-and-load router reach about 15\% of the ceiling. Battery-weighted Dijkstra reaches about 85\% with no training. LEST load-table Dijkstra reaches about 90\% with no solver. An LP planner re-solved every round reaches 97\%. LEST load-table Dijkstra builds each round's tree node by node, books every source's traffic on its path, and shares load through a 4-tier table whose overhead it pays for. It wins 118 of 120 time-varying deployments against battery Dijkstra, and the ranking holds on real Intel Lab traffic. Prediction adds half a point to the planner and nothing to Dijkstra, and learning adds no measurable gain. For lifetime routing in this setting, coordination and planning are what matter.

Future work will cover four directions:
\begin{itemize}[leftmargin=*]
  \item a distributed LEST Dijkstra, in which every node computes the same tree from a shared snapshot of energy and load;
  \item a packet-level evaluation with lossy links taken from the Intel connectivity data, collisions and retransmissions, together with networks of 200--500 nodes;
  \item further published learned baselines, including MADII retrained on time-varying traffic or the authors' code if it is released;
  \item traffic with stronger drift or weekly cycles, where prediction may matter more.
\end{itemize}
